% JAC1964-20261121, /root/arithmetic_sign.
\subsection*{AR-26. The full Hermitian divisor family does not orient the height}

This is a constraint symmetry, not an example of a negative geometric
height.  It tests the larger divisor family suggested by the CM action,
rather than only the one-parameter inequalities of AR-22.

\begin{proposition}[AR-26]
Let $A=E^3$, where $E$ has complex multiplication by $\mathbf Z[i]$,
and use the product principal polarization $L_0$.  Work over a number
field containing the required CM endomorphisms.  Let $f:C\to A$ be a
morphism from a smooth projective geometrically integral curve, with
a rational point $e$ satisfying $f(e)=0$, and suppose that
$f_*:\operatorname{Jac}(C)\to A$ is surjective.  Set
\[
 V=f_*[C],\qquad Z=\operatorname{Ce}(V)\in
 \operatorname{CH}^2(N)_\mathbf Q,
 \qquad N=h^3(A)_{(3,0)+(0,3)},\qquad q=h_K(Z,Z).
\]
The exact projector defining $N$ is the one in AR-25.
For $R\in\operatorname{Herm}_3(\mathbf Q(i))$, let
$\overline L_R$ denote the symmetric rigidified rational line bundle
with its canonical metric corresponding to $R$.  For
$B\in M_3(\mathbf Q(i))$, let $\overline{\mathcal Q}_B$ denote the
canonical cross bundle on $A\times A$ obtained from the Poincare
bundle for $L_0$ and $B$, and put
$\overline Q_B=(f\times f)^*\overline{\mathcal Q}_B$.
In particular $\mathcal Q_R=m^*L_R-p_1^*L_R-p_2^*L_R$ when $R$
is Hermitian.  Write
\[
 P_{1,R}=p_1^*f^*\overline L_R,\qquad
 P_{2,S}=p_2^*f^*\overline L_S.
\]
All arithmetic intersections below are divided by $[K:\mathbf Q]$.
Then the following assertions hold.
\begin{enumerate}
\item Every symmetric divisor is primitive on this motive:
\[
 L_R Z=0\quad\hbox{in }\operatorname{CH}^3(A)_\mathbf Q.
                                                               \tag{AR.26a}
\]
Every arithmetic triple with exactly one cross factor vanishes:
\[
 \langle P_{j,R}P_{k,S}Q_B\rangle_K=0
 \qquad(j,k\in\{1,2\}).                                \tag{AR.26b}
\]
The entire symmetric cross cubic is
\[
 \langle Q_R^3\rangle_K=-6\det(R)q
 \qquad(R\in\operatorname{Herm}_3(\mathbf Q(i))).         \tag{AR.26c}
\]
\item Let $\mathcal D$ be the rational vector space of the displayed
canonical adelic divisor classes on $C^2$.  There is a well-defined
linear involution
\[
 \sigma(P_{1,R}+P_{2,S}+Q_B)=P_{1,R}+P_{2,S}-Q_B.         \tag{AR.26d}
\]
It preserves the geometric intersection pairing $G$ on $\mathcal D$.
It also preserves the cone $\mathcal C$ of classes obtained by pulling
back canonical nef bundles associated with positive semidefinite
Hermitian matrices
\[
 H=\begin{pmatrix}R&B\\ B^\dagger&S\end{pmatrix}
 \quad\hbox{on }E^3\times E^3.
\]
\item Put $I(D_1,D_2,D_3)=\langle D_1D_2D_3\rangle_K$ and define
\[
 I^\sigma(D_1,D_2,D_3)=I(\sigma D_1,\sigma D_2,\sigma D_3).
\]
The functional $I^\sigma$ has exactly the same even-cross intersection
data as $I$, and still has the vanishing (AR.26b).  It changes the sign
of every three-cross intersection, in particular changing $q$ to
$-q$ in (AR.26c).

Every divisor Hodge-index inequality
\[
 I(M,M,P)\le0,\qquad
 P\in\mathcal C,\quad G(P,P)>0,\quad G(M,P)=0,           \tag{AR.26e}
\]
is preserved under $I\mapsto I^\sigma$.  So are the nonnegative
intersections $I(P_1,P_2,P_3)\ge0$ with all $P_j\in\mathcal C$.
Consequently this complete Hermitian family, including its polarized
inequalities and all the even intersection data, cannot distinguish
$q$ from $-q$ and does not orient a nonzero value.
\end{enumerate}
More concretely, on the ambient divisor labels $I^\sigma$ is the
actual intersection functional obtained from
$(f,-f):C^2\to A^2$, in place of $(f,f)$.  Its odd cycle factors are
$Z$ and $-Z$, so its changed coefficient represents the mixed height
$h_K(Z,-Z)=-q$.  It is not a negative self-height.  Thus this family
of inequalities does not distinguish the required equality of the
two curve maps from their opposites.  It gives no negative Chow
class and no second diagonal realization.  It also does not address
nef divisors on $C^2$ outside the stated pullback cone, or an
additional diagonal-sensitive arithmetic identity.
\end{proposition}

\begin{proof}
Let $a=[1+i]$ act diagonally on $A$ and put $U=a^*$.  On $N$ the
exact identity from AR-25 is
\[
 U^2+4U+8=0.                                             \tag{AR.26f}
\]
On $h^1(A)$ the endomorphism identity is $U^2-2U+2=0$.
Since $\deg(a)=8$, its pushforward on $h^1(A)$ is $8U^{-1}$ and
satisfies $T^2-8T+32=0$.  Poincare duality identifies this operator
with the dual of $a^*$ on $h^5(A)$.  Thus the exact Chow-motive
identity on $h^5(A)$ is
\[
 U^2-8U+32=0.                                            \tag{AR.26g}
\]
This uses identities of endomorphisms and Chow--Kunneth pieces,
not merely the characteristic polynomial on realization.

The symmetric divisor $L_R$ has Beauville grade zero, so multiplication
by it sends $h^3(A)$ to $h^5(A)$.  Since $a$ is a central CM
endomorphism, $a^*L_R=2L_R$.  Hence, writing this multiplication as
$\ell_R$, one has $U\ell_R=2\ell_RU$.  From (AR.26f) it follows that
\[
 (U^2+8U+32)\ell_R|_N
 =4\ell_R(U^2+4U+8)|_N=0.
\]
Subtract (AR.26g) to obtain $16U\ell_R|_N=0$.
The isogeny operator $U$ is invertible with rational coefficients,
and therefore $\ell_R|_N=0$.  This proves (AR.26a).

For (AR.26b), first note that the Albanese sum of $L_RV$ is zero.
Indeed symmetry of $L_R$ and the action of inversion on Albanese
give
\[
 \operatorname{alb}_A(L_RV)
 =\operatorname{alb}_A(L_R\operatorname{Ce}(V))
 =\operatorname{alb}_A(L_RZ)=0
 \quad\hbox{in }A(K)\otimes\mathbf Q.                   \tag{AR.26h}
\]
Thus the even part of $V$, including any other Beauville components,
has not been discarded without justification: its Albanese
contribution is zero by this identity.  The relative Deligne pairing
of $P_{1,R}$ and $Q_B$ along the first curve is the canonical
Picard line obtained by applying the Poincare correspondence $B$ to
the sum (AR.26h).  It is the zero class in the rational adelic Picard
group.  The metric assertion follows from its trivial rigidification
and its $[n]^*$-eigenvalue $n$, exactly as in AR-18; one may first
clear denominators and then kill the torsion Albanese point.
Pairing with $P_{2,S}$ proves the case $j\ne k$.
If $j=k$, the projection formula instead uses the degree-zero
restriction of $Q_B$ on the other curve and gives zero directly.

For positive definite $R$, Gu's height identity, in the normalization
of AR-10 with $m=2$, gives
\[
 \langle\mathcal Q_R^3\cdot(Z\times Z)\rangle_K
       =-6\chi(A,L_R)h_K(Z,Z).
\]
Here (AR.26a) supplies its Chow-primitivity hypothesis.  Relative to
the product principal polarization on $E^3$, the usual degree formula
is $\chi(A,L_R)=\det(R)$.  This also follows by diagonalizing the
Hermitian form over $\mathbf Q(i)$ and using the degree of the
corresponding isogeny; rational powers are handled by homogeneity.
The cross bundle is odd under inversion in each factor.  Therefore
its cube pairs with $V\times V$ exactly as with $Z\times Z$, by the
same representative-level identity as AR-18.  Projection along
$f\times f$ gives (AR.26c) for positive definite $R$.  Both sides
are cubic polynomials on the rational Hermitian space, so equality
on its open positive cone proves the identity for every $R$.

We next justify that (AR.26d) is an involution on the indicated
adelic divisor space, although $\sigma$ need not come from an
automorphism of $C^2$.  The cross bundles are metrically trivial
on both axes through $e$, and the vertical bundles are normalized
at $(e,e)$.  Restricting a relation among the three
summands to $C\times\{e\}$ and to $\{e\}\times C$ separates its
two vertical terms.  The remaining cross term is represented in
\[
 \operatorname{Pic}^{-}(C^2)_\mathbf Q
   =\operatorname{Hom}(\operatorname{Jac}(C),
                       \operatorname{Jac}(C)^\vee)_\mathbf Q
\]
by the composite induced by $f_*$, the Poincare map $B$, and $f^*$.
Surjectivity of $f_*$ and injectivity up to finite kernel of its dual
make this map injective in $B$ over $\mathbf Q$.  Thus that remaining
term vanishes only if $B=0$.  The admissible cross metrics are unique
with the prescribed rigidifications, so the same separation holds
for their adelic classes.  This proves that relations are
$\sigma$-stable, rather than postulating a surface automorphism.

The geometric product of one vertical divisor and one cross divisor
has degree zero.  Hence $G$ pairs the vertical and cross summands
orthogonally, which proves its $\sigma$-invariance.  On the ambient
Hermitian matrices, $\sigma$ is conjugation by
$\operatorname{diag}(I_3,-I_3)$.  It preserves positive
semidefiniteness.  Such matrices indeed give canonical nef rational
bundles: Hermitian elimination expresses a rational positive
semidefinite matrix as a sum of positive rational multiples of
rank-one forms over $\mathbf Q(i)$; these are pullbacks of the
canonical ample bundle on $E$, after clearing denominators.
Pullback to $C^2$ preserves nefness.  This proves preservation of
$\mathcal C$.  For these nef classes, bigness of the underlying
surface divisor is equivalent to $G(P,P)>0$, which is also
preserved.

The parity assertions about $I^\sigma$ now follow by multilinearity
and (AR.26b).  For (AR.26e), $\sigma M$ is still integrable,
$\sigma P$ is still nef with big underlying divisor, and
$G(\sigma M,\sigma P)=G(M,P)=0$.  Applying Yuan--Zhang's theorem
to this pair proves
\[
 I^\sigma(M,M,P)=I(\sigma M,\sigma M,\sigma P)\le0.
\]
The same argument applies to nonnegative intersections of three nef
canonical classes.  Therefore the whole system is invariant, not
only any fixed finite list of tests.  Passing to consequences by
linear combinations, polarization, or limits cannot break a
symmetry of that system.  In particular its unchanged even data
cannot select one sign for the determinant coefficient in
(AR.26c).

Finally, if $s(x,y)=(x,-y)$ on $A^2$, the canonical metric identities
are $s^*\overline L_1=\overline L_1$,
$s^*\overline L_2=\overline L_2$ and
$s^*\overline Q_B=-\overline Q_B$.  The first two use symmetry of
the rigidified bundles, and the third is Poincare bilinearity;
uniqueness of the canonical metrics makes them metric identities.
Consequently pullback by $(f,-f)=s\circ(f,f)$ realizes $I^\sigma$
on the ambient labels.  Its three-cross intersection involves
$Z\times(-Z)$, giving the bilinear height $-q$ by AR-10.  The
formal sign ambiguity therefore has a concrete mixed-height
realization, while providing no negative self-pairing and no
identification of that mixed map with a diagonal map.
\end{proof}

\paragraph{Sources and exact scope.}
The divisor inequality used here is Yuan--Zhang, \emph{The arithmetic
Hodge index theorem for adelic line bundles}, Math. Ann. 367 (2017),
Theorem 3.2, pp.~1134--1135: integrable $M$, nef $P$ with big
underlying bundle, and $M\cdot P=0$.  No equality criterion is used.
The canonical metric and normalized-correspondence assertions are
Zhang 2010, Lemmas 2.2.1--2.2.4, pp.~20--22, and his discussion of
composition of arithmetic correspondences on p.~23.  Gu's
Theorem 1.3 and Proposition 5.3 are used exactly through the already
checked AR-10 identity.  The new part is the all-divisor CM
primitivity, its determinant specialization, and the full cone
symmetry.  These sharpen the diagnosis of the failed divisor route;
the full complementary height sign remains open in this argument.
